% Options for packages loaded elsewhere
\PassOptionsToPackage{unicode}{hyperref}
\PassOptionsToPackage{hyphens}{url}
\PassOptionsToPackage{dvipsnames,svgnames,x11names}{xcolor}
%
\documentclass[
  11pt,
  a4paper,
]{ltjsarticle}
\usepackage{amsmath,amssymb}
\usepackage{iftex}
\ifPDFTeX
  \usepackage[T1]{fontenc}
  \usepackage[utf8]{inputenc}
  \usepackage{textcomp} % provide euro and other symbols
\else % if luatex or xetex
  \usepackage{unicode-math} % this also loads fontspec
  \defaultfontfeatures{Scale=MatchLowercase}
  \defaultfontfeatures[\rmfamily]{Ligatures=TeX,Scale=1}
\fi
\usepackage{lmodern}
\ifPDFTeX\else
  % xetex/luatex font selection
\fi
% Use upquote if available, for straight quotes in verbatim environments
\IfFileExists{upquote.sty}{\usepackage{upquote}}{}
\IfFileExists{microtype.sty}{% use microtype if available
  \usepackage[]{microtype}
  \UseMicrotypeSet[protrusion]{basicmath} % disable protrusion for tt fonts
}{}
\makeatletter
\@ifundefined{KOMAClassName}{% if non-KOMA class
  \IfFileExists{parskip.sty}{%
    \usepackage{parskip}
  }{% else
    \setlength{\parindent}{0pt}
    \setlength{\parskip}{6pt plus 2pt minus 1pt}}
}{% if KOMA class
  \KOMAoptions{parskip=half}}
\makeatother
\usepackage{xcolor}
\usepackage[margin=24mm]{geometry}
\setlength{\emergencystretch}{3em} % prevent overfull lines
\providecommand{\tightlist}{%
  \setlength{\itemsep}{0pt}\setlength{\parskip}{0pt}}
\setcounter{secnumdepth}{5}
\ifLuaTeX
  \usepackage{selnolig}  % disable illegal ligatures
\fi
\IfFileExists{bookmark.sty}{\usepackage{bookmark}}{\usepackage{hyperref}}
\IfFileExists{xurl.sty}{\usepackage{xurl}}{} % add URL line breaks if available
\urlstyle{same}
\hypersetup{
  colorlinks=true,
  linkcolor={blue},
  filecolor={Maroon},
  citecolor={Blue},
  urlcolor={blue},
  pdfcreator={LaTeX via pandoc}}

\author{}
\date{}

\begin{document}

\hypertarget{ux89e3ux7b54-03-ux9023ux7acbux4e00ux6b21ux65b9ux7a0bux5f0f}{%
\section{解答 03 ---
連立一次方程式}\label{ux89e3ux7b54-03-ux9023ux7acbux4e00ux6b21ux65b9ux7a0bux5f0f}}

\hypertarget{section}{%
\subsection{1}\label{section}}

第2式から第1式の2倍を引くと \texttt{3z=2}、第3式に第1式を加えると
\texttt{z=2} となり矛盾します。したがって解は存在しません。

この矛盾は augmented matrix に \texttt{0=非零}
の行が現れることに対応します。

\hypertarget{section-1}{%
\subsection{2}\label{section-1}}

\texttt{A={[}a\_1\ ...\ a\_n{]}} なら

\[
Ax=x_1a_1+\cdots+x_na_n.
\]

よってある \texttt{x} で \texttt{Ax=b} が成立することは、\texttt{b}
が列の一次結合で書けること、すなわち \texttt{b∈Col(A)} と同値です。

\hypertarget{section-2}{%
\subsection{3}\label{section-2}}

任意の解 \texttt{x} に対して \texttt{A(x-x\_p)=b-b=0} なので
\texttt{x-x\_p∈ker(A)}。逆に \texttt{z∈ker(A)} なら
\texttt{A(x\_p+z)=b}。よって全解は \texttt{x\_p+ker(A)} です。

\hypertarget{section-3}{%
\subsection{4}\label{section-3}}

核に非零 \texttt{z} があれば \texttt{x\_p} と \texttt{x\_p+z}
が別解になるため一意ではありません。逆に核が \texttt{\{0\}} なら二解
\texttt{x\_1,x\_2} に対し \texttt{A(x\_1-x\_2)=0} から
\texttt{x\_1=x\_2} です。

\hypertarget{section-4}{%
\subsection{5}\label{section-4}}

正方行列では full rank \texttt{n} なら nullity
は0です。したがって写像は単射かつ全射となり、逆写像が存在します。これが逆行列です。

\end{document}
